% AML (CSAI2017P) --- Theory Quiz 1 (T01) --- 8 sets, A--H, four pages each.
% GENERATED BY build.py --- edit content.py, not this file.
% Released after all sittings. Answer key: T01-solutions.pdf in this folder.
\documentclass[10pt,a4paper]{article}
\usepackage[margin=16mm,top=12mm,bottom=14mm]{geometry}
\usepackage[T1]{fontenc}
\usepackage{lmodern}
\usepackage{enumitem}
\usepackage{fancyvrb}
\usepackage{amsmath}
\usepackage{amssymb}
\usepackage{array}
\usepackage{fancyhdr}
\usepackage{microtype}

\setlength{\parindent}{0pt}
\setlength{\emergencystretch}{3em}
\newcommand{\BL}{\rule[-1.5pt]{22mm}{0.4pt}}
\newcommand{\blk}[1]{\rule[-1.5pt]{#1}{0.4pt}}

\pagestyle{fancy}\fancyhf{}
\renewcommand{\headrulewidth}{0.4pt}
\renewcommand{\footrulewidth}{0.4pt}
\newcommand{\SetLetter}{}
\fancyhead[L]{\footnotesize CSAI2017P --- Applied Machine Learning (Theory) --- Theory Quiz 1}
\fancyhead[R]{\footnotesize\bfseries Set \SetLetter}
\fancyfoot[L]{\footnotesize Name:\,\blk{34mm} \quad SAP ID:\,\blk{28mm}}
\fancyfoot[R]{\footnotesize Set \SetLetter\ --- page \thepage\ of 4}

\DefineVerbatimEnvironment{code}{Verbatim}%
  {fontsize=\small,xleftmargin=5mm,baselinestretch=0.98}

\newlist{qA}{enumerate}{1}
\setlist[qA]{label=\textbf{A\arabic*.},leftmargin=8mm,topsep=4pt,itemsep=9pt,parsep=1pt}
\newlist{qB}{enumerate}{1}
\setlist[qB]{label=\textbf{B\arabic*.},leftmargin=8mm,topsep=5pt,itemsep=13pt,parsep=1pt}
\newlist{qC}{enumerate}{1}
\setlist[qC]{label=\textbf{C\arabic*.},leftmargin=8mm,topsep=5pt,itemsep=8pt,parsep=1pt}
\newlist{qD}{enumerate}{1}
\setlist[qD]{label=\textbf{D\arabic*.},leftmargin=8mm,topsep=3pt,itemsep=4pt,parsep=1pt}
\newlist{qE}{enumerate}{1}
\setlist[qE]{label=\textbf{E\arabic*.},leftmargin=8mm,topsep=3pt,itemsep=5pt,parsep=1pt}

\newcommand{\parthead}[2]{\vspace{2pt}\noindent\rule{\linewidth}{0.4pt}\par\vspace{2pt}%
  \noindent{\bfseries Part #1}\hfill{\small #2}\par\vspace{1pt}}

\makeatletter
\newcommand{\rules}[2]{\par\vspace{2mm}%
  \@tempcnta=0
  \loop\ifnum\@tempcnta<#1\relax
    \noindent\rule{\linewidth}{0.3pt}\par\vspace{#2}%
    \advance\@tempcnta by 1
  \repeat}
\makeatother

\begin{document}


\renewcommand{\SetLetter}{A}\setcounter{page}{1}


\begin{center}
{\large\bfseries CSAI2017P --- Applied Machine Learning (Theory)}\\[2pt]
{\large\bfseries Theory Quiz 1 (T01) \quad\textbullet\quad Set A}\\[3pt]
{\small Duration \textbf{40 minutes} \quad\textbullet\quad Maximum marks \textbf{20}
 \quad\textbullet\quad 15 questions on four pages}
\end{center}
\vspace{2pt}

% Column widths differ slightly from LQ01's: at LQ01's widths the row runs
% ~11.6pt past the margin and "Signature" does not fit its own cell.
\noindent\begin{tabular}{@{}p{1.4cm}p{5.0cm}p{1.9cm}p{3.5cm}p{1.3cm}p{2.4cm}@{}}
\textbf{Name} & \hrulefill & \textbf{SAP ID} & \hrulefill & \textbf{Batch} & \hrulefill\\[9pt]
\textbf{Date} & \hrulefill & \textbf{Signature} & \hrulefill & \textbf{Set} & \fbox{\rule{0pt}{11pt}\ A\ }\\
\end{tabular}

\vspace{4pt}
\noindent\fbox{\parbox{\dimexpr\linewidth-2\fboxsep-2\fboxrule}{\small
\textbf{Syllabus.} Prerequisite material --- probability, Python, linear algebra and descriptive
statistics --- together with Unit I (Lecture 1): what machine learning is, the three families,
generalisation, cross-validation, data leakage and the Python stack; and Unit II (Lecture 2): loss
functions, MSE, MAE, RMSE, Huber, cross-entropy, hinge and triplet loss.\\[3pt]
\textbf{Instructions.} Closed book. Write your answers on this paper, in the space provided. A
calculator is not needed --- the arithmetic is deliberately short. Answer every question --- there
is no negative marking. In Part C a bare True/False without a reason scores half. In Parts D and E
the marks are for the reasoning, not for the vocabulary; an answer that states what it cannot
determine earns credit, an answer that asserts without a reason does not. Marks are shown against
each part.\\[3pt]
\textbf{About the set letter.} Every set of this quiz examines the same topics at the same level,
for the same marks, against the same marking scheme. Your set reflects the \emph{style} of question
that suits how you work, and nothing else. It has no effect on your mark.}}
\vspace{3pt}


\parthead{A --- Multiple choice}{5 questions $\times$ 1 = 5 marks. Circle one letter.}

\begin{qA}
  \item Let $A$ be the event ``at least one head'' in two tosses of a fair coin. Then $P(A)$ equals\\
    (a) $3/4$ \quad (b) $1/4$\\
    (c) $1/2$ \quad (d) $1$
  \item \texttt{len(range(2, 10, 3))} evaluates to\\
    (a) 2 \quad (b) 3\\
    (c) 4 \quad (d) an error
  \item $A$ is $4 \times 3$ and $B$ is $3 \times 5$. The product $AB$ has shape\\
    (a) $3 \times 3$ \quad (b) $5 \times 4$\\
    (c) $4 \times 5$ \quad (d) undefined
  \item A model is fitted to labelled pairs $(x_i, y_i)$ in which every $y_i$ is a real number on a continuous scale. This is\\
    (a) supervised classification \quad (b) unsupervised\\
    (c) reinforcement \quad (d) supervised regression
  \item Over constant predictions $c$, the quantity $\frac{1}{n}\sum_i \lvert y_i - c \rvert$ is minimised by the\\
    (a) median of the $y_i$ \quad (b) mean of the $y_i$\\
    (c) mode of the $y_i$ \quad (d) largest $y_i$
\end{qA}

\vspace{4pt}

\parthead{B --- Fill in the blanks}{4 $\times$ 1 = 4 marks. Write the missing word on the rule.}

\begin{qB}

  \item For the values $2,\,3,\,4,\,5,\,96$ the \BL\ exceeds the median.

  \item If $\mathrm{RMSE}_{\text{train}} \ll \mathrm{RMSE}_{\text{test}}$, the model is said to \BL\ the training data.

  \item If a step that learns from the data is fitted on the whole of $X$ before $X$ is split, test-row information reaches the model. This fault is called data \BL.

  \item For one example, cross-entropy is the negative \BL\ of the probability assigned to the correct class.

\end{qB}

\vfill
\begin{center}\small
\begin{tabular}{|>{\centering\arraybackslash}m{1.9cm}|*{5}{>{\centering\arraybackslash}m{1.6cm}|}}
\hline
\rule{0pt}{12pt}\textbf{Question} & A1 & A2 & A3 & A4 & A5\\[2pt]\hline
\rule{0pt}{20pt}\textbf{Answer} & & & & & \\[7pt]\hline
\end{tabular}\\[2pt]
{\footnotesize Copy your circled Part A letters into this grid. The grid is what is marked.}
\end{center}

\newpage

\parthead{C --- True or false}{3 $\times$ 1 = 3 marks. Circle one, then give a reason.}

\begin{qC}

  \item A model achieving zero training error is guaranteed to perform well on unseen data. \hfill \textbf{TRUE} \; / \; \textbf{FALSE}\\[1pt]
    {\small Reason:}\,\blk{140mm}\rules{1}{6.4mm}

  \item Model $P$ reports $\mathrm{MAE} = 1.0$ and model $Q$ reports $\mathrm{MSE} = 1.375$, both from the identical predictions; therefore $P$ is the better model. \hfill \textbf{TRUE} \; / \; \textbf{FALSE}\\[1pt]
    {\small Reason:}\,\blk{140mm}\rules{1}{6.4mm}

  \item For a point with $y = +1$ and score $f(x) = 3$, the hinge loss $\max(0,\,1 - y f(x))$ is exactly zero, so that point does not push the boundary during training. \hfill \textbf{TRUE} \; / \; \textbf{FALSE}\\[1pt]
    {\small Reason:}\,\blk{140mm}\rules{1}{6.4mm}

\end{qC}

\vspace{4pt}

\parthead{D --- Short answer}{2 questions $\times$ 2 = 4 marks. Show how you decide, not just what you decide.}

\begin{qD}

  \item For $y = (10,\,10,\,10)$ and $\hat{y} = (10,\,13,\,7)$: \textbf{(a)} compute the MSE and the MAE, showing your working; \textbf{(b)} state which of the two is more affected by the largest single residual, and why.\par\vspace{78mm}

\end{qD}

\newpage

\parthead{D --- Short answer, continued}{}

\begin{qD}[start=2]

  \item A five-fold cross-validation of one model returns the fold scores $(0.95,\ 0.93,\ 0.97,\ 0.62,\ 0.96)$. \textbf{(a)} Explain why reporting only their mean would be inadequate. \textbf{(b)} State what you would investigate first, and why.\par\vspace{80mm}

\end{qD}

\vspace{4pt}

\parthead{E --- Reasoning}{1 question $\times$ 4 = 4 marks. You are marked on the reasoning you make visible.}

\begin{qE}

  \item A colleague reports a test accuracy of $0.97$ from the procedure below, on a dataset with
$n = 100$ rows and $p = 5000$ columns. A domain expert says the honest figure should be close to chance.

\vspace{2mm}
\hspace*{6mm}\textbf{Step 1.}\ Load $X \in \mathbb{R}^{100 \times 5000}$ and $y \in \{0,1\}^{100}$.\\
\hspace*{6mm}\textbf{Step 2.}\ Score all $5000$ columns against $y$; keep the $20$ most strongly associated.\\
\hspace*{6mm}\textbf{Step 3.}\ Split the $100$ rows into a training part and a test part.\\
\hspace*{6mm}\textbf{Step 4.}\ Fit a classifier on the training part.\\
\hspace*{6mm}\textbf{Step 5.}\ Report accuracy on the test part.


    \vspace{2mm}
    Answer all three parts overleaf. Any diagnosis that fits the evidence and is argued for earns
    full credit --- naming the ``expected'' fault is not what is marked.
    \begin{enumerate}[label=(\roman*),leftmargin=7mm,itemsep=2pt,topsep=2pt]
      \item Name the most likely cause of the reported figure, and explain \emph{why} it inflates it. \hfill \textbf{[2]}
      \item State one concrete check that would confirm or rule out your diagnosis. \hfill \textbf{[1]}
      \item Say what you would expect to see if your diagnosis turned out to be wrong. \hfill \textbf{[1]}
    \end{enumerate}

\end{qE}

\newpage

\parthead{E --- Reasoning, answer space}{Label each part (i), (ii), (iii). A sketch, a table or a list is as acceptable as prose.}

\vfill

{\footnotesize\itshape Continue on a continuation sheet if you need more room; label anything you write there with its question number.}

\newpage

\renewcommand{\SetLetter}{B}\setcounter{page}{1}


\begin{center}
{\large\bfseries CSAI2017P --- Applied Machine Learning (Theory)}\\[2pt]
{\large\bfseries Theory Quiz 1 (T01) \quad\textbullet\quad Set B}\\[3pt]
{\small Duration \textbf{40 minutes} \quad\textbullet\quad Maximum marks \textbf{20}
 \quad\textbullet\quad 15 questions on four pages}
\end{center}
\vspace{2pt}

% Column widths differ slightly from LQ01's: at LQ01's widths the row runs
% ~11.6pt past the margin and "Signature" does not fit its own cell.
\noindent\begin{tabular}{@{}p{1.4cm}p{5.0cm}p{1.9cm}p{3.5cm}p{1.3cm}p{2.4cm}@{}}
\textbf{Name} & \hrulefill & \textbf{SAP ID} & \hrulefill & \textbf{Batch} & \hrulefill\\[9pt]
\textbf{Date} & \hrulefill & \textbf{Signature} & \hrulefill & \textbf{Set} & \fbox{\rule{0pt}{11pt}\ B\ }\\
\end{tabular}

\vspace{4pt}
\noindent\fbox{\parbox{\dimexpr\linewidth-2\fboxsep-2\fboxrule}{\small
\textbf{Syllabus.} Prerequisite material --- probability, Python, linear algebra and descriptive
statistics --- together with Unit I (Lecture 1): what machine learning is, the three families,
generalisation, cross-validation, data leakage and the Python stack; and Unit II (Lecture 2): loss
functions, MSE, MAE, RMSE, Huber, cross-entropy, hinge and triplet loss.\\[3pt]
\textbf{Instructions.} Closed book. Write your answers on this paper, in the space provided. A
calculator is not needed --- the arithmetic is deliberately short. Answer every question --- there
is no negative marking. In Part C a bare True/False without a reason scores half. In Parts D and E
the marks are for the reasoning, not for the vocabulary; an answer that states what it cannot
determine earns credit, an answer that asserts without a reason does not. Marks are shown against
each part.\\[3pt]
\textbf{About the set letter.} Every set of this quiz examines the same topics at the same level,
for the same marks, against the same marking scheme. Your set reflects the \emph{style} of question
that suits how you work, and nothing else. It has no effect on your mark.}}
\vspace{3pt}


\parthead{A --- Multiple choice}{5 questions $\times$ 1 = 5 marks. Circle one letter.}

\begin{qA}
  \item The four equally likely outcomes of two fair coin tosses are HH, HT, TH, TT. The proportion containing at least one head is\\
    (a) $1/4$ \quad (b) $3/4$\\
    (c) $1/2$ \quad (d) $1$
  \item For \texttt{lst = [10, 20, 30, 40]}, the expression \texttt{len(lst[1:3])} evaluates to\\
    (a) 1 \quad (b) 3\\
    (c) 2 \quad (d) 4
  \item An array of shape \texttt{(4, 3)} is matrix-multiplied by an array of shape \texttt{(3, 5)}. The result has shape\\
    (a) \texttt{(3, 3)} \quad (b) \texttt{(5, 4)}\\
    (c) undefined \quad (d) \texttt{(4, 5)}
  \item A table has the columns \emph{floor area}, \emph{bedrooms} and \emph{sale price in lakh}. Predicting the third from the first two is\\
    (a) supervised regression \quad (b) supervised classification\\
    (c) unsupervised \quad (d) reinforcement
  \item A model's four residuals are $-1,\ +1,\ +1,\ +12$. The loss least disturbed by the value $+12$ is\\
    (a) MSE \quad (b) RMSE\\
    (c) MAE \quad (d) all equally
\end{qA}

\vspace{4pt}

\parthead{B --- Fill in the blanks}{4 $\times$ 1 = 4 marks. Write the missing word on the rule.}

\begin{qB}

  \item For the readings $2,\,3,\,4,\,5,\,96$ the mean is $22$ and the median is $4$; the \BL\ is the summary dragged upward by the single large reading.

  \item A model records a training RMSE of $0.4$ and a test RMSE of $3.9$. It is said to \BL\ the training data.

  \item A step that learns from the data is fitted on all $500$ rows before the $400/100$ split, so the test rows contribute to what it learned. This fault is called data \BL.

  \item A model gives probability $0.7$ to the correct class of one example; that example's cross-entropy is the negative \BL\ of $0.7$.

\end{qB}

\vfill
\begin{center}\small
\begin{tabular}{|>{\centering\arraybackslash}m{1.9cm}|*{5}{>{\centering\arraybackslash}m{1.6cm}|}}
\hline
\rule{0pt}{12pt}\textbf{Question} & A1 & A2 & A3 & A4 & A5\\[2pt]\hline
\rule{0pt}{20pt}\textbf{Answer} & & & & & \\[7pt]\hline
\end{tabular}\\[2pt]
{\footnotesize Copy your circled Part A letters into this grid. The grid is what is marked.}
\end{center}

\newpage

\parthead{C --- True or false}{3 $\times$ 1 = 3 marks. Circle one, then give a reason.}

\begin{qC}

  \item A results table showing a training accuracy of $1.00$ is sufficient evidence that the model will do well on new data. \hfill \textbf{TRUE} \; / \; \textbf{FALSE}\\[1pt]
    {\small Reason:}\,\blk{140mm}\rules{1}{6.4mm}

  \item A table reports $\mathrm{MAE} = 1.0$ for model $P$ and $\mathrm{MSE} = 1.375$ for model $Q$, computed from the identical predictions. $P$ is therefore the better model. \hfill \textbf{TRUE} \; / \; \textbf{FALSE}\\[1pt]
    {\small Reason:}\,\blk{140mm}\rules{1}{6.4mm}

  \item A table row records $y = +1$, $f(x) = 3$ and hinge loss $\max(0,\,1 - y f(x)) = 0$; that zero means the point does not push the boundary during training. \hfill \textbf{TRUE} \; / \; \textbf{FALSE}\\[1pt]
    {\small Reason:}\,\blk{140mm}\rules{1}{6.4mm}

\end{qC}

\vspace{4pt}

\parthead{D --- Short answer}{2 questions $\times$ 2 = 4 marks. Show how you decide, not just what you decide.}

\begin{qD}

  \item A table records three true values and the three predictions made for them: $y = (6,\,6,\,6)$ and $\hat{y} = (5,\,7,\,10)$. \textbf{(a)} Compute the MSE and the MAE, showing your working. \textbf{(b)} State which of the two is more affected by the largest single residual, and why.\par\vspace{78mm}

\end{qD}

\newpage

\parthead{D --- Short answer, continued}{}

\begin{qD}[start=2]

  \item A five-fold cross-validation of one model returns the fold scores $(0.88,\ 0.91,\ 0.90,\ 0.55,\ 0.89)$. \textbf{(a)} Explain why reporting only their mean would be inadequate. \textbf{(b)} State what you would investigate first, and why.\par\vspace{80mm}

\end{qD}

\vspace{4pt}

\parthead{E --- Reasoning}{1 question $\times$ 4 = 4 marks. You are marked on the reasoning you make visible.}

\begin{qE}

  \item The table below records a colleague's experiment. The domain expert consulted afterwards says the
honest accuracy should be close to chance.

\vspace{2mm}
\begin{center}\small
\begin{tabular}{@{}ll@{}}
\hline
\rule{0pt}{11pt}Rows in the dataset & $100$\\
Columns in the dataset & $5000$\\
Columns kept, scored against $y$ over \emph{all} $100$ rows & $20$\\
Split performed & \emph{after} those columns were kept\\
Reported test accuracy & $\mathbf{0.97}$\\[2pt]
\hline
\end{tabular}
\end{center}


    \vspace{2mm}
    Answer all three parts overleaf. Any diagnosis that fits the evidence and is argued for earns
    full credit --- naming the ``expected'' fault is not what is marked.
    \begin{enumerate}[label=(\roman*),leftmargin=7mm,itemsep=2pt,topsep=2pt]
      \item Name the most likely cause of the reported figure, and explain \emph{why} it inflates it. \hfill \textbf{[2]}
      \item State one concrete check that would confirm or rule out your diagnosis. \hfill \textbf{[1]}
      \item Say what you would expect to see if your diagnosis turned out to be wrong. \hfill \textbf{[1]}
    \end{enumerate}

\end{qE}

\newpage

\parthead{E --- Reasoning, answer space}{Label each part (i), (ii), (iii). A sketch, a table or a list is as acceptable as prose.}

\vfill

{\footnotesize\itshape Continue on a continuation sheet if you need more room; label anything you write there with its question number.}

\newpage

\renewcommand{\SetLetter}{C}\setcounter{page}{1}


\begin{center}
{\large\bfseries CSAI2017P --- Applied Machine Learning (Theory)}\\[2pt]
{\large\bfseries Theory Quiz 1 (T01) \quad\textbullet\quad Set C}\\[3pt]
{\small Duration \textbf{40 minutes} \quad\textbullet\quad Maximum marks \textbf{20}
 \quad\textbullet\quad 15 questions on four pages}
\end{center}
\vspace{2pt}

% Column widths differ slightly from LQ01's: at LQ01's widths the row runs
% ~11.6pt past the margin and "Signature" does not fit its own cell.
\noindent\begin{tabular}{@{}p{1.4cm}p{5.0cm}p{1.9cm}p{3.5cm}p{1.3cm}p{2.4cm}@{}}
\textbf{Name} & \hrulefill & \textbf{SAP ID} & \hrulefill & \textbf{Batch} & \hrulefill\\[9pt]
\textbf{Date} & \hrulefill & \textbf{Signature} & \hrulefill & \textbf{Set} & \fbox{\rule{0pt}{11pt}\ C\ }\\
\end{tabular}

\vspace{4pt}
\noindent\fbox{\parbox{\dimexpr\linewidth-2\fboxsep-2\fboxrule}{\small
\textbf{Syllabus.} Prerequisite material --- probability, Python, linear algebra and descriptive
statistics --- together with Unit I (Lecture 1): what machine learning is, the three families,
generalisation, cross-validation, data leakage and the Python stack; and Unit II (Lecture 2): loss
functions, MSE, MAE, RMSE, Huber, cross-entropy, hinge and triplet loss.\\[3pt]
\textbf{Instructions.} Closed book. Write your answers on this paper, in the space provided. A
calculator is not needed --- the arithmetic is deliberately short. Answer every question --- there
is no negative marking. In Part C a bare True/False without a reason scores half. In Parts D and E
the marks are for the reasoning, not for the vocabulary; an answer that states what it cannot
determine earns credit, an answer that asserts without a reason does not. Marks are shown against
each part.\\[3pt]
\textbf{About the set letter.} Every set of this quiz examines the same topics at the same level,
for the same marks, against the same marking scheme. Your set reflects the \emph{style} of question
that suits how you work, and nothing else. It has no effect on your mark.}}
\vspace{3pt}


\parthead{A --- Multiple choice}{5 questions $\times$ 1 = 5 marks. Circle one letter.}

\begin{qA}
  \item A probability tree for two tosses of a fair coin has four leaves of equal weight. The combined weight of the leaves showing at least one head is\\
    (a) $1/4$ \quad (b) $1/2$\\
    (c) $3/4$ \quad (d) $1$
  \item \texttt{len(set([1, 2, 2, 3]))} evaluates to\\
    (a) 1 \quad (b) 2\\
    (c) 4 \quad (d) 3
  \item One rectangle is drawn $4$ rows tall and $3$ columns wide, a second $3$ rows tall and $5$ columns wide. Multiplied in that order, the result rectangle is\\
    (a) $4$ by $5$ \quad (b) $3$ by $3$\\
    (c) $5$ by $4$ \quad (d) not defined
  \item A scatter plot of two features carries no colour coding at all, and the task is to draw three natural clumps around the points. This is\\
    (a) supervised regression \quad (b) unsupervised\\
    (c) supervised classification \quad (d) reinforcement
  \item Two loss curves are plotted against the error $e$. One is a smooth parabola through the origin; the other is a V with a sharp corner at the origin. The V is\\
    (a) MSE \quad (b) hinge loss\\
    (c) cross-entropy \quad (d) MAE
\end{qA}

\vspace{4pt}

\parthead{B --- Fill in the blanks}{4 $\times$ 1 = 4 marks. Write the missing word on the rule.}

\begin{qB}

  \item A histogram of $2,\,3,\,4,\,5,\,96$ has a long right tail. The vertical line marking the \BL\ sits far to the right of the line marking the median.

  \item On a learning curve the training error keeps falling while the validation error turns upward. From that turning point onward the model is said to \BL.

  \item A pipeline diagram shows the box that learns from the data placed \emph{before} the split box rather than inside the fold. The fault this creates is called data \BL.

  \item A plot of $-\log p$ against $p$ rises without bound as $p \to 0$. Cross-entropy for one example is exactly this curve: the negative \BL\ of the probability given to the correct class.

\end{qB}

\vfill
\begin{center}\small
\begin{tabular}{|>{\centering\arraybackslash}m{1.9cm}|*{5}{>{\centering\arraybackslash}m{1.6cm}|}}
\hline
\rule{0pt}{12pt}\textbf{Question} & A1 & A2 & A3 & A4 & A5\\[2pt]\hline
\rule{0pt}{20pt}\textbf{Answer} & & & & & \\[7pt]\hline
\end{tabular}\\[2pt]
{\footnotesize Copy your circled Part A letters into this grid. The grid is what is marked.}
\end{center}

\newpage

\parthead{C --- True or false}{3 $\times$ 1 = 3 marks. Circle one, then give a reason.}

\begin{qC}

  \item A learning curve whose training-error line reaches zero shows a model that will generalise well. \hfill \textbf{TRUE} \; / \; \textbf{FALSE}\\[1pt]
    {\small Reason:}\,\blk{140mm}\rules{1}{6.4mm}

  \item Two bars on one chart, labelled $\mathrm{MAE} = 1.0$ for model $P$ and $\mathrm{MSE} = 1.375$ for model $Q$ and computed from identical predictions, show that $P$ is the better model. \hfill \textbf{TRUE} \; / \; \textbf{FALSE}\\[1pt]
    {\small Reason:}\,\blk{140mm}\rules{1}{6.4mm}

  \item On a plot of hinge loss against the margin $y f(x)$, the curve is flat at zero for every margin beyond $1$; a point sitting at margin $3$ therefore exerts no pull on the boundary. \hfill \textbf{TRUE} \; / \; \textbf{FALSE}\\[1pt]
    {\small Reason:}\,\blk{140mm}\rules{1}{6.4mm}

\end{qC}

\vspace{4pt}

\parthead{D --- Short answer}{2 questions $\times$ 2 = 4 marks. Show how you decide, not just what you decide.}

\begin{qD}

  \item Three points are plotted against a horizontal line at $y = 8$. Their predicted heights are $\hat{y} = (6,\,10,\,13)$ while every true value is $y = 8$. \textbf{(a)} Compute the MSE and the MAE, showing your working. \textbf{(b)} State which of the two is more affected by the largest single residual, and why.\par\vspace{78mm}

\end{qD}

\newpage

\parthead{D --- Short answer, continued}{}

\begin{qD}[start=2]

  \item A bar chart of five cross-validation fold scores shows four bars of similar height and one much shorter: $(0.92,\ 0.90,\ 0.94,\ 0.61,\ 0.93)$. \textbf{(a)} Explain why reporting only their mean would be inadequate. \textbf{(b)} State what you would investigate first, and why.\par\vspace{80mm}

\end{qD}

\vspace{4pt}

\parthead{E --- Reasoning}{1 question $\times$ 4 = 4 marks. You are marked on the reasoning you make visible.}

\begin{qE}

  \item A colleague's plot shows the reported test accuracy sitting at $0.97$, \emph{above} every one of the
training-fold accuracies drawn beside it --- an ordering that should not occur. The dataset is $100$ rows by
$5000$ columns. Reading the pipeline diagram from left to right, the boxes are

\vspace{2mm}
\begin{center}\small
\fbox{\ load $X$ ($100 \times 5000$), $y$\ } $\rightarrow$
\fbox{\ keep the $20$ columns best associated with $y$\ } $\rightarrow$
\fbox{\ split\ } $\rightarrow$
\fbox{\ fit\ } $\rightarrow$
\fbox{\ score\ }
\end{center}


    \vspace{2mm}
    Answer all three parts overleaf. Any diagnosis that fits the evidence and is argued for earns
    full credit --- naming the ``expected'' fault is not what is marked.
    \begin{enumerate}[label=(\roman*),leftmargin=7mm,itemsep=2pt,topsep=2pt]
      \item Name the most likely cause of the reported figure, and explain \emph{why} it inflates it. \hfill \textbf{[2]}
      \item State one concrete check that would confirm or rule out your diagnosis. \hfill \textbf{[1]}
      \item Say what you would expect to see if your diagnosis turned out to be wrong. \hfill \textbf{[1]}
    \end{enumerate}

\end{qE}

\newpage

\parthead{E --- Reasoning, answer space}{Label each part (i), (ii), (iii). A sketch, a table or a list is as acceptable as prose.}

\vfill

{\footnotesize\itshape Continue on a continuation sheet if you need more room; label anything you write there with its question number.}

\newpage

\renewcommand{\SetLetter}{D}\setcounter{page}{1}


\begin{center}
{\large\bfseries CSAI2017P --- Applied Machine Learning (Theory)}\\[2pt]
{\large\bfseries Theory Quiz 1 (T01) \quad\textbullet\quad Set D}\\[3pt]
{\small Duration \textbf{40 minutes} \quad\textbullet\quad Maximum marks \textbf{20}
 \quad\textbullet\quad 15 questions on four pages}
\end{center}
\vspace{2pt}

% Column widths differ slightly from LQ01's: at LQ01's widths the row runs
% ~11.6pt past the margin and "Signature" does not fit its own cell.
\noindent\begin{tabular}{@{}p{1.4cm}p{5.0cm}p{1.9cm}p{3.5cm}p{1.3cm}p{2.4cm}@{}}
\textbf{Name} & \hrulefill & \textbf{SAP ID} & \hrulefill & \textbf{Batch} & \hrulefill\\[9pt]
\textbf{Date} & \hrulefill & \textbf{Signature} & \hrulefill & \textbf{Set} & \fbox{\rule{0pt}{11pt}\ D\ }\\
\end{tabular}

\vspace{4pt}
\noindent\fbox{\parbox{\dimexpr\linewidth-2\fboxsep-2\fboxrule}{\small
\textbf{Syllabus.} Prerequisite material --- probability, Python, linear algebra and descriptive
statistics --- together with Unit I (Lecture 1): what machine learning is, the three families,
generalisation, cross-validation, data leakage and the Python stack; and Unit II (Lecture 2): loss
functions, MSE, MAE, RMSE, Huber, cross-entropy, hinge and triplet loss.\\[3pt]
\textbf{Instructions.} Closed book. Write your answers on this paper, in the space provided. A
calculator is not needed --- the arithmetic is deliberately short. Answer every question --- there
is no negative marking. In Part C a bare True/False without a reason scores half. In Parts D and E
the marks are for the reasoning, not for the vocabulary; an answer that states what it cannot
determine earns credit, an answer that asserts without a reason does not. Marks are shown against
each part.\\[3pt]
\textbf{About the set letter.} Every set of this quiz examines the same topics at the same level,
for the same marks, against the same marking scheme. Your set reflects the \emph{style} of question
that suits how you work, and nothing else. It has no effect on your mark.}}
\vspace{3pt}


\parthead{A --- Multiple choice}{5 questions $\times$ 1 = 5 marks. Circle one letter.}

\begin{qA}
  \item Two independent inspections of a component each pass with probability $1/2$. The probability that at least one passes is\\
    (a) $1/4$ \quad (b) $1/2$\\
    (c) $1$ \quad (d) $3/4$
  \item A shop's stock is written as \texttt{\{'red': 1, 'blue': 2, 'red': 3\}}. Its \texttt{len} is\\
    (a) 2 \quad (b) 1\\
    (c) 3 \quad (d) an error
  \item A clinic records $8$ measurements for each of $200$ patients. To score a patient with one weight per measurement, the weight vector must have length\\
    (a) 200 \quad (b) 8\\
    (c) 1600 \quad (d) any length
  \item A bank holds five years of loans, each marked \emph{repaid} or \emph{defaulted}, and wants a verdict for a new applicant. This is\\
    (a) supervised regression \quad (b) unsupervised\\
    (c) supervised classification \quad (d) reinforcement
  \item A temperature sensor occasionally reports impossible values that cannot be removed before training. The loss whose fitted constant is least disturbed by them is\\
    (a) MSE \quad (b) MAE\\
    (c) RMSE \quad (d) squared hinge
\end{qA}

\vspace{4pt}

\parthead{B --- Fill in the blanks}{4 $\times$ 1 = 4 marks. Write the missing word on the rule.}

\begin{qB}

  \item A ward records stays of $2$, $3$, $4$, $5$ and $96$ days. The \BL\ length of stay is the one pulled upward by the single long stay.

  \item A churn model is right about almost every past customer and wrong about most new ones. It is said to \BL\ the training data.

  \item An analyst lets a step that learns from the data see the full customer table, then splits it into train and test parts. The fault is called data \BL.

  \item For one email the spam filter gives probability $0.7$ to the class the email actually belongs to. That email's cross-entropy is the negative \BL\ of $0.7$.

\end{qB}

\vfill
\begin{center}\small
\begin{tabular}{|>{\centering\arraybackslash}m{1.9cm}|*{5}{>{\centering\arraybackslash}m{1.6cm}|}}
\hline
\rule{0pt}{12pt}\textbf{Question} & A1 & A2 & A3 & A4 & A5\\[2pt]\hline
\rule{0pt}{20pt}\textbf{Answer} & & & & & \\[7pt]\hline
\end{tabular}\\[2pt]
{\footnotesize Copy your circled Part A letters into this grid. The grid is what is marked.}
\end{center}

\newpage

\parthead{C --- True or false}{3 $\times$ 1 = 3 marks. Circle one, then give a reason.}

\begin{qC}

  \item A fraud model that flags every past fraud in the bank's own records correctly is ready to be deployed on next month's transactions. \hfill \textbf{TRUE} \; / \; \textbf{FALSE}\\[1pt]
    {\small Reason:}\,\blk{140mm}\rules{1}{6.4mm}

  \item A vendor's model reports $\mathrm{MAE} = 1.0$ and yours reports $\mathrm{MSE} = 1.375$ on the identical predictions; the vendor's model is therefore better. \hfill \textbf{TRUE} \; / \; \textbf{FALSE}\\[1pt]
    {\small Reason:}\,\blk{140mm}\rules{1}{6.4mm}

  \item A support-vector classifier scores a correctly labelled application at $f(x) = 3$ with $y = +1$, so its hinge loss is zero and removing that application from the training set would not move the boundary. \hfill \textbf{TRUE} \; / \; \textbf{FALSE}\\[1pt]
    {\small Reason:}\,\blk{140mm}\rules{1}{6.4mm}

\end{qC}

\vspace{4pt}

\parthead{D --- Short answer}{2 questions $\times$ 2 = 4 marks. Show how you decide, not just what you decide.}

\begin{qD}

  \item A delivery service quotes $20$ minutes for three consecutive orders, so $y = (20,\,20,\,20)$. The model's predictions were $\hat{y} = (19,\,24,\,16)$. \textbf{(a)} Compute the MSE and the MAE, showing your working. \textbf{(b)} State which of the two is more affected by the largest single residual, and why.\par\vspace{78mm}

\end{qD}

\newpage

\parthead{D --- Short answer, continued}{}

\begin{qD}[start=2]

  \item A five-fold cross-validation of a churn model returns the fold scores $(0.86,\ 0.84,\ 0.88,\ 0.52,\ 0.87)$. \textbf{(a)} Explain why reporting only their mean would be inadequate. \textbf{(b)} State what you would investigate first, and why.\par\vspace{80mm}

\end{qD}

\vspace{4pt}

\parthead{E --- Reasoning}{1 question $\times$ 4 = 4 marks. You are marked on the reasoning you make visible.}

\begin{qE}

  \item An intern presents a churn model at the weekly meeting and reports a test accuracy of $0.97$. The
dataset holds $100$ customers and $5000$ recorded attributes per customer. The team's own domain expert says
the honest figure should be close to chance, because most of those attributes are meaningless. Asked how the
result was produced, the intern says

\vspace{2mm}
\hspace*{6mm}\parbox{0.9\linewidth}{\emph{``I looked at all $5000$ attributes across all $100$ customers,
kept the $20$ that tracked churn most closely, then split the customers into a training set and a test set,
trained on the training set, and scored on the test set.''}}


    \vspace{2mm}
    Answer all three parts overleaf. Any diagnosis that fits the evidence and is argued for earns
    full credit --- naming the ``expected'' fault is not what is marked.
    \begin{enumerate}[label=(\roman*),leftmargin=7mm,itemsep=2pt,topsep=2pt]
      \item Name the most likely cause of the reported figure, and explain \emph{why} it inflates it. \hfill \textbf{[2]}
      \item State one concrete check that would confirm or rule out your diagnosis. \hfill \textbf{[1]}
      \item Say what you would expect to see if your diagnosis turned out to be wrong. \hfill \textbf{[1]}
    \end{enumerate}

\end{qE}

\newpage

\parthead{E --- Reasoning, answer space}{Label each part (i), (ii), (iii). A sketch, a table or a list is as acceptable as prose.}

\vfill

{\footnotesize\itshape Continue on a continuation sheet if you need more room; label anything you write there with its question number.}

\newpage

\renewcommand{\SetLetter}{E}\setcounter{page}{1}


\begin{center}
{\large\bfseries CSAI2017P --- Applied Machine Learning (Theory)}\\[2pt]
{\large\bfseries Theory Quiz 1 (T01) \quad\textbullet\quad Set E}\\[3pt]
{\small Duration \textbf{40 minutes} \quad\textbullet\quad Maximum marks \textbf{20}
 \quad\textbullet\quad 15 questions on four pages}
\end{center}
\vspace{2pt}

% Column widths differ slightly from LQ01's: at LQ01's widths the row runs
% ~11.6pt past the margin and "Signature" does not fit its own cell.
\noindent\begin{tabular}{@{}p{1.4cm}p{5.0cm}p{1.9cm}p{3.5cm}p{1.3cm}p{2.4cm}@{}}
\textbf{Name} & \hrulefill & \textbf{SAP ID} & \hrulefill & \textbf{Batch} & \hrulefill\\[9pt]
\textbf{Date} & \hrulefill & \textbf{Signature} & \hrulefill & \textbf{Set} & \fbox{\rule{0pt}{11pt}\ E\ }\\
\end{tabular}

\vspace{4pt}
\noindent\fbox{\parbox{\dimexpr\linewidth-2\fboxsep-2\fboxrule}{\small
\textbf{Syllabus.} Prerequisite material --- probability, Python, linear algebra and descriptive
statistics --- together with Unit I (Lecture 1): what machine learning is, the three families,
generalisation, cross-validation, data leakage and the Python stack; and Unit II (Lecture 2): loss
functions, MSE, MAE, RMSE, Huber, cross-entropy, hinge and triplet loss.\\[3pt]
\textbf{Instructions.} Closed book. Write your answers on this paper, in the space provided. A
calculator is not needed --- the arithmetic is deliberately short. Answer every question --- there
is no negative marking. In Part C a bare True/False without a reason scores half. In Parts D and E
the marks are for the reasoning, not for the vocabulary; an answer that states what it cannot
determine earns credit, an answer that asserts without a reason does not. Marks are shown against
each part.\\[3pt]
\textbf{About the set letter.} Every set of this quiz examines the same topics at the same level,
for the same marks, against the same marking scheme. Your set reflects the \emph{style} of question
that suits how you work, and nothing else. It has no effect on your mark.}}
\vspace{3pt}


\parthead{A --- Multiple choice}{5 questions $\times$ 1 = 5 marks. Circle one letter.}

\begin{qA}
  \item After \texttt{o = list(product('HT', repeat=2))}, the value of \texttt{sum(1 for x in o if 'H' in x) / len(o)} is\\
    (a) 0.75 \quad (b) 0.25\\
    (c) 0.5 \quad (d) 1.0
  \item After \texttt{xs = [i for i in range(5) if i \% 2 == 0]}, \texttt{len(xs)} is\\
    (a) 2 \quad (b) 4\\
    (c) 3 \quad (d) 5
  \item \texttt{np.zeros((4, 3)) @ np.ones((3, 5))} has shape\\
    (a) \texttt{(3, 3)} \quad (b) \texttt{(4, 5)}\\
    (c) \texttt{(5, 4)} \quad (d) it raises
  \item \texttt{KMeans(n\_clusters=3).fit(X)} is called and no \texttt{y} is passed anywhere. The task is\\
    (a) supervised regression \quad (b) supervised classification\\
    (c) reinforcement \quad (d) unsupervised
  \item \texttt{np.mean(np.abs(y - y\_hat))} computes\\
    (a) MSE \quad (b) RMSE\\
    (c) MAE \quad (d) cross-entropy
\end{qA}

\vspace{4pt}

\parthead{B --- Fill in the blanks}{4 $\times$ 1 = 4 marks. Write the missing word on the rule.}

\begin{qB}

  \item \texttt{np.mean([2, 3, 4, 5, 96])} returns $22.0$ and \texttt{np.median(...)} returns $4.0$. The \BL\ is the one moved by the value $96$.

  \item \texttt{model.score(X\_train, y\_train)} returns $0.99$ and \texttt{model.score(X\_test, y\_test)} returns $0.61$. The model is said to \BL\ the training data.

  \item A step that learns from the data is called on the whole of \texttt{X} before \texttt{train\_test\_split} is reached. The fault this creates is called data \BL.

  \item For one example, \texttt{-np.log(p\_correct)} is the cross-entropy: the negative \BL\ of the probability given to the correct class.

\end{qB}

\vfill
\begin{center}\small
\begin{tabular}{|>{\centering\arraybackslash}m{1.9cm}|*{5}{>{\centering\arraybackslash}m{1.6cm}|}}
\hline
\rule{0pt}{12pt}\textbf{Question} & A1 & A2 & A3 & A4 & A5\\[2pt]\hline
\rule{0pt}{20pt}\textbf{Answer} & & & & & \\[7pt]\hline
\end{tabular}\\[2pt]
{\footnotesize Copy your circled Part A letters into this grid. The grid is what is marked.}
\end{center}

\newpage

\parthead{C --- True or false}{3 $\times$ 1 = 3 marks. Circle one, then give a reason.}

\begin{qC}

  \item \texttt{model.score(X\_train, y\_train) == 1.0} is enough to conclude that the model will do well on unseen data. \hfill \textbf{TRUE} \; / \; \textbf{FALSE}\\[1pt]
    {\small Reason:}\,\blk{140mm}\rules{1}{6.4mm}

  \item From the identical arrays, \texttt{mean\_absolute\_error(...)} returns $1.0$ for model $P$ and \texttt{mean\_squared\_error(...)} returns $1.375$ for model $Q$. Since $1.0 < 1.375$, $P$ is the better model. \hfill \textbf{TRUE} \; / \; \textbf{FALSE}\\[1pt]
    {\small Reason:}\,\blk{140mm}\rules{1}{6.4mm}

  \item For \texttt{y = 1} and \texttt{f = 3.0}, \texttt{max(0.0, 1 - y * f)} evaluates to \texttt{0.0}, so this point contributes nothing to the hinge objective and does not move the boundary. \hfill \textbf{TRUE} \; / \; \textbf{FALSE}\\[1pt]
    {\small Reason:}\,\blk{140mm}\rules{1}{6.4mm}

\end{qC}

\vspace{4pt}

\parthead{D --- Short answer}{2 questions $\times$ 2 = 4 marks. Show how you decide, not just what you decide.}

\begin{qD}

  \item \texttt{y = np.array([12, 12, 12])} and \texttt{y\_hat = np.array([14, 10, 20])}. \textbf{(a)} Compute the MSE and the MAE by hand, showing your working. \textbf{(b)} State which of the two is more affected by the largest single residual, and why.\par\vspace{78mm}

\end{qD}

\newpage

\parthead{D --- Short answer, continued}{}

\begin{qD}[start=2]

  \item \texttt{cross\_val\_score(model, X, y, cv=5)} returns \texttt{array([0.94, 0.92, 0.96, 0.60, 0.95])}. \textbf{(a)} Explain why reporting only \texttt{.mean()} would be inadequate. \textbf{(b)} State what you would investigate first, and why.\par\vspace{80mm}

\end{qD}

\vspace{4pt}

\parthead{E --- Reasoning}{1 question $\times$ 4 = 4 marks. You are marked on the reasoning you make visible.}

\begin{qE}

  \item The function below is run on \texttt{X} of shape \texttt{(100, 5000)} and a binary \texttt{y}, and
returns $0.97$. A domain expert says the honest figure should be close to chance.

\begin{code}
from sklearn.feature_selection import SelectKBest, f_classif
from sklearn.model_selection import train_test_split
from sklearn.neighbors import KNeighborsClassifier

def evaluate(X, y):
    X_sel = SelectKBest(f_classif, k=20).fit_transform(X, y)
    X_tr, X_te, y_tr, y_te = train_test_split(X_sel, y, test_size=0.25,
                                              random_state=0)
    model = KNeighborsClassifier(n_neighbors=3).fit(X_tr, y_tr)
    return model.score(X_te, y_te)
\end{code}


    \vspace{2mm}
    Answer all three parts overleaf. Any diagnosis that fits the evidence and is argued for earns
    full credit --- naming the ``expected'' fault is not what is marked.
    \begin{enumerate}[label=(\roman*),leftmargin=7mm,itemsep=2pt,topsep=2pt]
      \item Name the most likely cause of the reported figure, and explain \emph{why} it inflates it. \hfill \textbf{[2]}
      \item State one concrete check that would confirm or rule out your diagnosis. \hfill \textbf{[1]}
      \item Say what you would expect to see if your diagnosis turned out to be wrong. \hfill \textbf{[1]}
    \end{enumerate}

\end{qE}

\newpage

\parthead{E --- Reasoning, answer space}{Label each part (i), (ii), (iii). A sketch, a table or a list is as acceptable as prose.}

\vfill

{\footnotesize\itshape Continue on a continuation sheet if you need more room; label anything you write there with its question number.}

\newpage

\renewcommand{\SetLetter}{F}\setcounter{page}{1}


\begin{center}
{\large\bfseries CSAI2017P --- Applied Machine Learning (Theory)}\\[2pt]
{\large\bfseries Theory Quiz 1 (T01) \quad\textbullet\quad Set F}\\[3pt]
{\small Duration \textbf{40 minutes} \quad\textbullet\quad Maximum marks \textbf{20}
 \quad\textbullet\quad 15 questions on four pages}
\end{center}
\vspace{2pt}

% Column widths differ slightly from LQ01's: at LQ01's widths the row runs
% ~11.6pt past the margin and "Signature" does not fit its own cell.
\noindent\begin{tabular}{@{}p{1.4cm}p{5.0cm}p{1.9cm}p{3.5cm}p{1.3cm}p{2.4cm}@{}}
\textbf{Name} & \hrulefill & \textbf{SAP ID} & \hrulefill & \textbf{Batch} & \hrulefill\\[9pt]
\textbf{Date} & \hrulefill & \textbf{Signature} & \hrulefill & \textbf{Set} & \fbox{\rule{0pt}{11pt}\ F\ }\\
\end{tabular}

\vspace{4pt}
\noindent\fbox{\parbox{\dimexpr\linewidth-2\fboxsep-2\fboxrule}{\small
\textbf{Syllabus.} Prerequisite material --- probability, Python, linear algebra and descriptive
statistics --- together with Unit I (Lecture 1): what machine learning is, the three families,
generalisation, cross-validation, data leakage and the Python stack; and Unit II (Lecture 2): loss
functions, MSE, MAE, RMSE, Huber, cross-entropy, hinge and triplet loss.\\[3pt]
\textbf{Instructions.} Closed book. Write your answers on this paper, in the space provided. A
calculator is not needed --- the arithmetic is deliberately short. Answer every question --- there
is no negative marking. In Part C a bare True/False without a reason scores half. In Parts D and E
the marks are for the reasoning, not for the vocabulary; an answer that states what it cannot
determine earns credit, an answer that asserts without a reason does not. Marks are shown against
each part.\\[3pt]
\textbf{About the set letter.} Every set of this quiz examines the same topics at the same level,
for the same marks, against the same marking scheme. Your set reflects the \emph{style} of question
that suits how you work, and nothing else. It has no effect on your mark.}}
\vspace{3pt}


\parthead{A --- Multiple choice}{5 questions $\times$ 1 = 5 marks. Circle one letter.}

\begin{qA}
  \item A fair coin is tossed twice. The chance of getting at least one head is\\
    (a) one in four \quad (b) three in four\\
    (c) one in two \quad (d) certain
  \item In Python, splitting the word \emph{machine} on the letter \emph{h} produces a list. Its length is\\
    (a) 1 \quad (b) 3\\
    (c) 7 \quad (d) 2
  \item A rectangular table of numbers with four rows and three columns is multiplied by a second table with three rows and five columns. The result has\\
    (a) three rows, three columns \quad (b) five rows, four columns\\
    (c) four rows, five columns \quad (d) no defined result
  \item A newsroom wants its articles grouped by topic, but nobody has ever labelled an article with a topic. This is\\
    (a) unsupervised \quad (b) supervised regression\\
    (c) supervised classification \quad (d) reinforcement
  \item Two ways of scoring the same errors: one squares each error before averaging, the other takes the size of each before averaging. The second is less disturbed by a single wild reading. It is\\
    (a) MSE \quad (b) RMSE\\
    (c) hinge loss \quad (d) MAE
\end{qA}

\vspace{4pt}

\parthead{B --- Fill in the blanks}{4 $\times$ 1 = 4 marks. Write the missing word on the rule.}

\begin{qB}

  \item When one value in a small set of readings is far larger than the rest, the \BL\ is pulled towards it while the median stays put.

  \item A model that has learned the particular examples it was shown rather than the pattern behind them is said to \BL.

  \item Using, at training time, information that would not have been available at prediction time is called data \BL.

  \item For a single example, cross-entropy is the negative \BL\ of the probability the model gave to the correct class.

\end{qB}

\vfill
\begin{center}\small
\begin{tabular}{|>{\centering\arraybackslash}m{1.9cm}|*{5}{>{\centering\arraybackslash}m{1.6cm}|}}
\hline
\rule{0pt}{12pt}\textbf{Question} & A1 & A2 & A3 & A4 & A5\\[2pt]\hline
\rule{0pt}{20pt}\textbf{Answer} & & & & & \\[7pt]\hline
\end{tabular}\\[2pt]
{\footnotesize Copy your circled Part A letters into this grid. The grid is what is marked.}
\end{center}

\newpage

\parthead{C --- True or false}{3 $\times$ 1 = 3 marks. Circle one, then give a reason.}

\begin{qC}

  \item Getting every training example right proves that a model has learned the underlying pattern. \hfill \textbf{TRUE} \; / \; \textbf{FALSE}\\[1pt]
    {\small Reason:}\,\blk{140mm}\rules{1}{6.4mm}

  \item One model's mean absolute error is a smaller number than another model's mean squared error on the same predictions, so the first model is the better one. \hfill \textbf{TRUE} \; / \; \textbf{FALSE}\\[1pt]
    {\small Reason:}\,\blk{140mm}\rules{1}{6.4mm}

  \item A correctly classified point that already sits comfortably beyond the margin costs a hinge-loss model nothing, and therefore has no say in where the boundary is placed. \hfill \textbf{TRUE} \; / \; \textbf{FALSE}\\[1pt]
    {\small Reason:}\,\blk{140mm}\rules{1}{6.4mm}

\end{qC}

\vspace{4pt}

\parthead{D --- Short answer}{2 questions $\times$ 2 = 4 marks. Show how you decide, not just what you decide.}

\begin{qD}

  \item Three quantities were each truly fifteen, and the model predicted twelve, eighteen and twenty-one for them in that order. \textbf{(a)} Compute the mean squared error and the mean absolute error, showing your working. \textbf{(b)} State which of the two is more affected by the largest single residual, and why.\par\vspace{78mm}

\end{qD}

\newpage

\parthead{D --- Short answer, continued}{}

\begin{qD}[start=2]

  \item Splitting the data five ways and scoring each part in turn gave five numbers: $0.90$, $0.88$, $0.92$, $0.57$ and $0.91$. \textbf{(a)} Explain why reporting only their average would be inadequate. \textbf{(b)} State what you would investigate first, and why.\par\vspace{80mm}

\end{qD}

\vspace{4pt}

\parthead{E --- Reasoning}{1 question $\times$ 4 = 4 marks. You are marked on the reasoning you make visible.}

\begin{qE}

  \item A colleague describes what she did, and reports an accuracy of ninety-seven per cent. The data she
used holds one hundred records, with five thousand recorded quantities per record. A domain expert says the
honest figure should be close to chance.

\vspace{2mm}
\hspace*{6mm}\parbox{0.9\linewidth}{\emph{``I went through all five thousand quantities, using every one of
the hundred records, and kept the twenty that lined up best with the answer I was trying to predict. Then I
set part of the data aside as a test set, trained on the rest, and measured how well it did on the part I had
set aside.''}}


    \vspace{2mm}
    Answer all three parts overleaf. Any diagnosis that fits the evidence and is argued for earns
    full credit --- naming the ``expected'' fault is not what is marked.
    \begin{enumerate}[label=(\roman*),leftmargin=7mm,itemsep=2pt,topsep=2pt]
      \item Name the most likely cause of the reported figure, and explain \emph{why} it inflates it. \hfill \textbf{[2]}
      \item State one concrete check that would confirm or rule out your diagnosis. \hfill \textbf{[1]}
      \item Say what you would expect to see if your diagnosis turned out to be wrong. \hfill \textbf{[1]}
    \end{enumerate}

\end{qE}

\newpage

\parthead{E --- Reasoning, answer space}{Label each part (i), (ii), (iii). A sketch, a table or a list is as acceptable as prose.}

\vfill

{\footnotesize\itshape Continue on a continuation sheet if you need more room; label anything you write there with its question number.}

\newpage

\renewcommand{\SetLetter}{G}\setcounter{page}{1}


\begin{center}
{\large\bfseries CSAI2017P --- Applied Machine Learning (Theory)}\\[2pt]
{\large\bfseries Theory Quiz 1 (T01) \quad\textbullet\quad Set G}\\[3pt]
{\small Duration \textbf{40 minutes} \quad\textbullet\quad Maximum marks \textbf{20}
 \quad\textbullet\quad 15 questions on four pages}
\end{center}
\vspace{2pt}

% Column widths differ slightly from LQ01's: at LQ01's widths the row runs
% ~11.6pt past the margin and "Signature" does not fit its own cell.
\noindent\begin{tabular}{@{}p{1.4cm}p{5.0cm}p{1.9cm}p{3.5cm}p{1.3cm}p{2.4cm}@{}}
\textbf{Name} & \hrulefill & \textbf{SAP ID} & \hrulefill & \textbf{Batch} & \hrulefill\\[9pt]
\textbf{Date} & \hrulefill & \textbf{Signature} & \hrulefill & \textbf{Set} & \fbox{\rule{0pt}{11pt}\ G\ }\\
\end{tabular}

\vspace{4pt}
\noindent\fbox{\parbox{\dimexpr\linewidth-2\fboxsep-2\fboxrule}{\small
\textbf{Syllabus.} Prerequisite material --- probability, Python, linear algebra and descriptive
statistics --- together with Unit I (Lecture 1): what machine learning is, the three families,
generalisation, cross-validation, data leakage and the Python stack; and Unit II (Lecture 2): loss
functions, MSE, MAE, RMSE, Huber, cross-entropy, hinge and triplet loss.\\[3pt]
\textbf{Instructions.} Closed book. Write your answers on this paper, in the space provided. A
calculator is not needed --- the arithmetic is deliberately short. Answer every question --- there
is no negative marking. In Part C a bare True/False without a reason scores half. In Parts D and E
the marks are for the reasoning, not for the vocabulary; an answer that states what it cannot
determine earns credit, an answer that asserts without a reason does not. Marks are shown against
each part.\\[3pt]
\textbf{About the set letter.} Every set of this quiz examines the same topics at the same level,
for the same marks, against the same marking scheme. Your set reflects the \emph{style} of question
that suits how you work, and nothing else. It has no effect on your mark.}}
\vspace{3pt}


\parthead{A --- Multiple choice}{5 questions $\times$ 1 = 5 marks. Circle one letter.}

\begin{qA}
  \item For two tosses of a fair coin, compare $P(\text{at least one head})$ with $P(\text{exactly one head})$\\
    (a) the second is larger \quad (b) they are equal\\
    (c) the first is larger \quad (d) cannot be decided
  \item Compare \texttt{len([1, 2, 3][:2])} with \texttt{len([1, 2] + [3])}\\
    (a) the second is larger \quad (b) the first is larger\\
    (c) they are equal \quad (d) one raises an error
  \item $A$ is $4 \times 3$ and $B$ is $3 \times 5$. Comparing the two products\\
    (a) only $BA$ is defined \quad (b) both, same shape\\
    (c) both, different shapes \quad (d) only $AB$ is defined
  \item Compare (i) predicting a house's sale price with (ii) deciding whether an email is spam\\
    (a) (i) classification, (ii) regression \quad (b) (i) regression, (ii) classification\\
    (c) both regression \quad (d) both unsupervised
  \item Compare MSE and MAE on a set of residuals containing one very large value. The one moved further by that value is\\
    (a) MSE \quad (b) MAE\\
    (c) neither, they move equally \quad (d) it depends on the sign
\end{qA}

\vspace{4pt}

\parthead{B --- Fill in the blanks}{4 $\times$ 1 = 4 marks. Write the missing word on the rule.}

\begin{qB}

  \item Compare the mean and median of $2,\,3,\,4,\,5,\,96$; the larger of the two is the \BL.

  \item Compare a model's error on the data it was fitted to with its error on unseen data. When the second is far worse, the model is said to \BL.

  \item Compare fitting a step that learns from the data \emph{inside} each fold with fitting it once on the whole dataset. The second choice produces data \BL.

  \item Compare accuracy with cross-entropy: accuracy merely counts, whereas cross-entropy is the negative \BL\ of the probability given to the correct class.

\end{qB}

\vfill
\begin{center}\small
\begin{tabular}{|>{\centering\arraybackslash}m{1.9cm}|*{5}{>{\centering\arraybackslash}m{1.6cm}|}}
\hline
\rule{0pt}{12pt}\textbf{Question} & A1 & A2 & A3 & A4 & A5\\[2pt]\hline
\rule{0pt}{20pt}\textbf{Answer} & & & & & \\[7pt]\hline
\end{tabular}\\[2pt]
{\footnotesize Copy your circled Part A letters into this grid. The grid is what is marked.}
\end{center}

\newpage

\parthead{C --- True or false}{3 $\times$ 1 = 3 marks. Circle one, then give a reason.}

\begin{qC}

  \item Between two models, the one with the higher training accuracy is the one that will do better on new data. \hfill \textbf{TRUE} \; / \; \textbf{FALSE}\\[1pt]
    {\small Reason:}\,\blk{140mm}\rules{1}{6.4mm}

  \item Between model $P$ ($\mathrm{MAE} = 1.0$) and model $Q$ ($\mathrm{MSE} = 1.375$) on identical predictions, $P$ is the better model. \hfill \textbf{TRUE} \; / \; \textbf{FALSE}\\[1pt]
    {\small Reason:}\,\blk{140mm}\rules{1}{6.4mm}

  \item Between a point at margin $y f(x) = 3$ and a point at margin $y f(x) = 0.5$, only the second contributes to the hinge objective; the first costs zero and does not move the boundary. \hfill \textbf{TRUE} \; / \; \textbf{FALSE}\\[1pt]
    {\small Reason:}\,\blk{140mm}\rules{1}{6.4mm}

\end{qC}

\vspace{4pt}

\parthead{D --- Short answer}{2 questions $\times$ 2 = 4 marks. Show how you decide, not just what you decide.}

\begin{qD}

  \item Compare three predictions $\hat{y} = (9,\,12,\,3)$ against three true values $y = (9,\,9,\,9)$. \textbf{(a)} Compute the MSE and the MAE, showing your working. \textbf{(b)} State which of the two is more affected by the largest single residual, and why.\par\vspace{78mm}

\end{qD}

\newpage

\parthead{D --- Short answer, continued}{}

\begin{qD}[start=2]

  \item Compare the five fold scores returned by a five-fold cross-validation: $(0.93,\ 0.91,\ 0.95,\ 0.59,\ 0.94)$. \textbf{(a)} Explain why reporting only their mean would be inadequate. \textbf{(b)} State what you would investigate first, and why.\par\vspace{80mm}

\end{qD}

\vspace{4pt}

\parthead{E --- Reasoning}{1 question $\times$ 4 = 4 marks. You are marked on the reasoning you make visible.}

\begin{qE}

  \item Two colleagues run the same dataset --- $100$ rows, $5000$ columns --- through the procedures below.
The procedures use the identical four steps in a different order, and report very different figures. The
domain expert says the honest figure should be close to chance.

\vspace{2mm}
\begin{center}\small
\begin{tabular}{@{}l@{\qquad}l@{}}
\hline
\rule{0pt}{11pt}\textbf{Procedure A --- reports $0.97$} & \textbf{Procedure B --- reports $0.51$}\\[2pt]
\hline
\rule{0pt}{11pt}1. Keep the $20$ columns best associated with $y$. & 1. Split into train and test.\\
2. Split into train and test. & 2. Keep the $20$ columns best associated with $y$.\\
3. Fit a classifier. & 3. Fit a classifier.\\
4. Report accuracy on the test part. & 4. Report accuracy on the test part.\\[2pt]
\hline
\end{tabular}
\end{center}


    \vspace{2mm}
    Answer all three parts overleaf. Any diagnosis that fits the evidence and is argued for earns
    full credit --- naming the ``expected'' fault is not what is marked.
    \begin{enumerate}[label=(\roman*),leftmargin=7mm,itemsep=2pt,topsep=2pt]
      \item Name the most likely cause of the reported figure, and explain \emph{why} it inflates it. \hfill \textbf{[2]}
      \item State one concrete check that would confirm or rule out your diagnosis. \hfill \textbf{[1]}
      \item Say what you would expect to see if your diagnosis turned out to be wrong. \hfill \textbf{[1]}
    \end{enumerate}

\end{qE}

\newpage

\parthead{E --- Reasoning, answer space}{Label each part (i), (ii), (iii). A sketch, a table or a list is as acceptable as prose.}

\vfill

{\footnotesize\itshape Continue on a continuation sheet if you need more room; label anything you write there with its question number.}

\newpage

\renewcommand{\SetLetter}{H}\setcounter{page}{1}


\begin{center}
{\large\bfseries CSAI2017P --- Applied Machine Learning (Theory)}\\[2pt]
{\large\bfseries Theory Quiz 1 (T01) \quad\textbullet\quad Set H}\\[3pt]
{\small Duration \textbf{40 minutes} \quad\textbullet\quad Maximum marks \textbf{20}
 \quad\textbullet\quad 15 questions on four pages}
\end{center}
\vspace{2pt}

% Column widths differ slightly from LQ01's: at LQ01's widths the row runs
% ~11.6pt past the margin and "Signature" does not fit its own cell.
\noindent\begin{tabular}{@{}p{1.4cm}p{5.0cm}p{1.9cm}p{3.5cm}p{1.3cm}p{2.4cm}@{}}
\textbf{Name} & \hrulefill & \textbf{SAP ID} & \hrulefill & \textbf{Batch} & \hrulefill\\[9pt]
\textbf{Date} & \hrulefill & \textbf{Signature} & \hrulefill & \textbf{Set} & \fbox{\rule{0pt}{11pt}\ H\ }\\
\end{tabular}

\vspace{4pt}
\noindent\fbox{\parbox{\dimexpr\linewidth-2\fboxsep-2\fboxrule}{\small
\textbf{Syllabus.} Prerequisite material --- probability, Python, linear algebra and descriptive
statistics --- together with Unit I (Lecture 1): what machine learning is, the three families,
generalisation, cross-validation, data leakage and the Python stack; and Unit II (Lecture 2): loss
functions, MSE, MAE, RMSE, Huber, cross-entropy, hinge and triplet loss.\\[3pt]
\textbf{Instructions.} Closed book. Write your answers on this paper, in the space provided. A
calculator is not needed --- the arithmetic is deliberately short. Answer every question --- there
is no negative marking. In Part C a bare True/False without a reason scores half. In Parts D and E
the marks are for the reasoning, not for the vocabulary; an answer that states what it cannot
determine earns credit, an answer that asserts without a reason does not. Marks are shown against
each part.\\[3pt]
\textbf{About the set letter.} Every set of this quiz examines the same topics at the same level,
for the same marks, against the same marking scheme. Your set reflects the \emph{style} of question
that suits how you work, and nothing else. It has no effect on your mark.}}
\vspace{3pt}


\parthead{A --- Multiple choice}{5 questions $\times$ 1 = 5 marks. Circle one letter.}

\begin{qA}
  \item \emph{For one toss of a fair coin, $P(\text{head}) = 1/2$.} For two tosses, $P(\text{at least one head})$ is\\
    (a) $1/4$ \quad (b) $1/2$\\
    (c) $1$ \quad (d) $3/4$
  \item \emph{\texttt{len([1, 2, 3][0:2])} is $2$.} Then \texttt{len([5, 6, 7, 8][1:4])} is\\
    (a) 2 \quad (b) 3\\
    (c) 4 \quad (d) an error
  \item \emph{A $2 \times 3$ times a $3 \times 4$ gives a $2 \times 4$.} Then a $4 \times 3$ times a $3 \times 5$ gives\\
    (a) $4 \times 5$ \quad (b) $3 \times 3$\\
    (c) $5 \times 4$ \quad (d) undefined
  \item \emph{Predicting a continuous quantity from labelled examples is supervised regression.} Sorting unlabelled photographs into groups nobody has named is\\
    (a) supervised regression \quad (b) supervised classification\\
    (c) unsupervised \quad (d) reinforcement
  \item \emph{The constant minimising $\sum_i (y_i - c)^2$ is the mean.} The constant minimising $\sum_i \lvert y_i - c \rvert$ is the\\
    (a) mean \quad (b) median\\
    (c) mode \quad (d) largest value
\end{qA}

\vspace{4pt}

\parthead{B --- Fill in the blanks}{4 $\times$ 1 = 4 marks. Write the missing word on the rule.}

\begin{qB}

  \item \emph{For $1,\,2,\,3$ the mean is $2$ and so is the median.} For $2,\,3,\,4,\,5,\,96$ they part company, and the larger of the two is the \BL.

  \item \emph{A model with training error $0.40$ and test error $0.45$ is generalising.} A model with training error $0.40$ and test error $3.90$ is instead said to \BL.

  \item \emph{A step that learns from the data, fitted inside each fold, is safe.} The same step fitted on the whole dataset before the split is data \BL.

  \item \emph{The probability given to the correct class was $0.9$, so the cross-entropy was $-\log 0.9$.} In general it is the negative \BL\ of that probability.

\end{qB}

\vfill
\begin{center}\small
\begin{tabular}{|>{\centering\arraybackslash}m{1.9cm}|*{5}{>{\centering\arraybackslash}m{1.6cm}|}}
\hline
\rule{0pt}{12pt}\textbf{Question} & A1 & A2 & A3 & A4 & A5\\[2pt]\hline
\rule{0pt}{20pt}\textbf{Answer} & & & & & \\[7pt]\hline
\end{tabular}\\[2pt]
{\footnotesize Copy your circled Part A letters into this grid. The grid is what is marked.}
\end{center}

\newpage

\parthead{C --- True or false}{3 $\times$ 1 = 3 marks. Circle one, then give a reason.}

\begin{qC}

  \item \emph{A model with zero training error and zero test error has generalised.} Therefore a model with zero training error has generalised. \hfill \textbf{TRUE} \; / \; \textbf{FALSE}\\[1pt]
    {\small Reason:}\,\blk{140mm}\rules{1}{6.4mm}

  \item \emph{Comparing two models by the same metric on the same held-out data is valid.} Comparing model $P$'s $\mathrm{MAE}$ of $1.0$ with model $Q$'s $\mathrm{MSE}$ of $1.375$ is therefore also valid, and shows $P$ to be better. \hfill \textbf{TRUE} \; / \; \textbf{FALSE}\\[1pt]
    {\small Reason:}\,\blk{140mm}\rules{1}{6.4mm}

  \item \emph{For $y = +1$ and $f(x) = 0.4$ the hinge loss is $0.6$, and the point pulls on the boundary.} For $y = +1$ and $f(x) = 3$ the hinge loss is $0$, and the point does not pull on the boundary. \hfill \textbf{TRUE} \; / \; \textbf{FALSE}\\[1pt]
    {\small Reason:}\,\blk{140mm}\rules{1}{6.4mm}

\end{qC}

\vspace{4pt}

\parthead{D --- Short answer}{2 questions $\times$ 2 = 4 marks. Show how you decide, not just what you decide.}

\begin{qD}

  \item \emph{For $y = (4,\,4)$ and $\hat{y} = (3,\,5)$ the residuals are $-1$ and $+1$, so $\mathrm{MSE} = 1$ and $\mathrm{MAE} = 1$.} Now take $y = (30,\,30,\,30)$ and $\hat{y} = (28,\,35,\,25)$. \textbf{(a)} Compute the MSE and the MAE, showing your working. \textbf{(b)} State which of the two is more affected by the largest single residual, and why.\par\vspace{78mm}

\end{qD}

\newpage

\parthead{D --- Short answer, continued}{}

\begin{qD}[start=2]

  \item \emph{Fold scores $(0.90,\ 0.91,\ 0.89)$ are tightly grouped, and their mean is a fair summary.} A five-fold run now returns $(0.89,\ 0.87,\ 0.91,\ 0.54,\ 0.90)$. \textbf{(a)} Explain why reporting only their mean would be inadequate. \textbf{(b)} State what you would investigate first, and why.\par\vspace{80mm}

\end{qD}

\vspace{4pt}

\parthead{E --- Reasoning}{1 question $\times$ 4 = 4 marks. You are marked on the reasoning you make visible.}

\begin{qE}

  \item \emph{Worked analogue. A student tunes a threshold by trying twenty values and reporting the best test
score. The reported figure is optimistic, because the test set has been consulted twenty times; the check is
to re-run with the tuning done inside the training data only, and the score should fall.}

\vspace{2mm}
Now consider this case. A colleague reports a test accuracy of $0.97$ on a dataset of $100$ rows and $5000$
columns, where the domain expert says the honest figure should be close to chance. The colleague scored all
$5000$ columns against $y$ using all $100$ rows, kept the best $20$, then split into train and test, fitted a
classifier and reported the test accuracy.


    \vspace{2mm}
    Answer all three parts overleaf. Any diagnosis that fits the evidence and is argued for earns
    full credit --- naming the ``expected'' fault is not what is marked.
    \begin{enumerate}[label=(\roman*),leftmargin=7mm,itemsep=2pt,topsep=2pt]
      \item Name the most likely cause of the reported figure, and explain \emph{why} it inflates it. \hfill \textbf{[2]}
      \item State one concrete check that would confirm or rule out your diagnosis. \hfill \textbf{[1]}
      \item Say what you would expect to see if your diagnosis turned out to be wrong. \hfill \textbf{[1]}
    \end{enumerate}

\end{qE}

\newpage

\parthead{E --- Reasoning, answer space}{Label each part (i), (ii), (iii). A sketch, a table or a list is as acceptable as prose.}

\vfill

{\footnotesize\itshape Continue on a continuation sheet if you need more room; label anything you write there with its question number.}

\newpage

\end{document}